% TOP_PROBLEM: {"id":30007670,"problem_number":"LOCAL-30007670","title":"Democracy's economic mechanisms and resilience","category_id":21,"category":{"id":21,"name":"economics","display_name":"Economics","description":"Open economic theory statements, published research agendas, and proposed scoped specifications; question provenance and historical priority are qualified per record.","slug":"economics","order_index":21,"created_at":"2026-10-08T00:00:00Z"},"status":"provenance_pending","external_url":null,"legacy_ids":[],"published":false,"statement":"In the explicitly specified setting: Which democratic institutions generate lasting growth and broad welfare gains, and under what economic conditions do those gains sustain support for democracy? The mathematical statement above fixes the target, assumptions and scope of this catalogue formulation.","economics":{"original_id":159,"field":"Political economy and institutions","kind":"identification","formalization_basis":"proposed_specification","source_status":"not_verified","priority_status":"not_established","priority_note":"No qualifying problem statement was directly verified, so historical priority cannot be assigned.","earliest_verified_reference":null,"current_reference":{"authors":["Daron Acemoglu","Nicolás Ajzenman","Cevat Giray Aksoy","Martin Fiszbein","Carlos A. Molina"],"title":"(Successful) Democracies Breed Their Own Support","year":2021,"url":"https://www.nber.org/papers/w29167","doi":"10.3386/w29167","locator":"Abstract, paragraphs on democracy exposure, economic performance, and support","evidence_summary":"The study identifies important performance-support links; the entire proposed growth-support-institutional-durability feedback question is not explicitly stated as remaining open."},"formalization_note":"This is a proposed scoped research specification. No primary passage explicitly stating this entire remaining question as open was verified. The supporting literature may answer important subquestions; this entry must not be advertised as a verified formal open problem."},"statusQualification":"Provenance pending: an explicit published statement of this exact open question was not verified. This is a catalogue-authored candidate specification, not a certified open conjecture. This is a proposed scoped research specification. No primary passage explicitly stating this entire remaining question as open was verified. The supporting literature may answer important subquestions; this entry must not be advertised as a verified formal open problem.","statusReviewedAt":"2026-10-08"}
% FIRST_POSTED: 2026-10-08
% ENTRY_KIND: statement_only
% !TeX program = lualatex
\documentclass[11pt]{article}
\usepackage{fontspec}
\usepackage[a4paper,margin=25mm]{geometry}
\usepackage{amsmath,amssymb,mathtools}
\usepackage{xurl}
\usepackage[hidelinks]{hyperref}
\newcommand{\sref}[2]{\hyperlink{source:#1:#2}{[S#2]}}
\setlength{\emergencystretch}{3em}
\begin{document}
% problemId:30007670
\section*{Democracy's economic mechanisms and resilience}\label{30007670}
\subsection{Definitions and mathematical statement}
Fix countries undergoing a specified transition from nondemocratic to democratic political rules, defining both regimes before analysis. Let \(D_{ct}\) indicate democratic status and \(Y_{cT}(d)\) real income per person under a declared regime path \(d\). Let \(S_{icT}(d)\) denote citizen \(i\)'s support for democratic procedures, measured independently of approval of the current incumbent. Define \(\tau_Y(T)=E[Y_{cT}(d^1)-Y_{cT}(d^0)]\), where \(d^1\) and \(d^0\) specify transition and comparator paths, and analogously define support effects. Determine which policy and institutional mechanisms mediate economic changes, and whether improved performance causally raises procedural support and institutional durability. Neither the average growth effect of a transition nor a correlation between income and support identifies this feedback loop. Specify assumptions for transition timing, anticipation, country selection, and global shocks. Additional randomized information or economic interventions may distinguish actual performance from beliefs about performance, but exclusion restrictions must be defended. Existing papers estimate important parts of this chain. The remaining combined mechanisms-and-resilience question is our proposed specification; a source explicitly declaring its full formulation open or identifying its first statement was not verified.

\par\textbf{Formulation and scope.} This is a proposed scoped research specification. No primary passage explicitly stating this entire remaining question as open was verified. The supporting literature may answer important subquestions; this entry must not be advertised as a verified formal open problem.

\par\textbf{Open-question provenance.} An explicit published open-question statement was not verified. Sources below are supporting literature and must not be treated as a first listing.

\par\textbf{Historical priority.} No qualifying problem statement was directly verified, so historical priority cannot be assigned.

\subsection{Short English statement}
In the explicitly specified setting: Which democratic institutions generate lasting growth and broad welfare gains, and under what economic conditions do those gains sustain support for democracy? The mathematical statement above fixes the target, assumptions and scope of this catalogue formulation.

\subsection{Sources}
\begin{itemize}\raggedright
\item[\textnormal{[S1]}]\hypertarget{source:30007670:1}{} Daron Acemoglu, Nicolás Ajzenman, Cevat Giray Aksoy, Martin Fiszbein, Carlos A. Molina. 2021. (Successful) Democracies Breed Their Own Support. Abstract, paragraphs on democracy exposure, economic performance, and support.
\url{https://www.nber.org/papers/w29167}.
DOI: 10.3386/w29167.
{\small\textit{Supporting or current literature; see evidence qualification. The study identifies important performance-support links; the entire proposed growth-support-institutional-durability feedback question is not explicitly stated as remaining open.}}
\end{itemize}
\end{document}
